\documentclass[../main.tex]{subfiles}

\begin{document}

\ifSubfilesClassLoaded{

    %\tableofcontents
    
    \setcounter{chapter}{4}
}{}

\chapter{ Week 5 }
 代靖涵 25120222201319
\begin{exercise}{}{}
设 $(A_n)$ 是拓扑空间 $X$ 的连通子集序列, 满足 $A_n \cap A_{n+1} \neq \emptyset, n \in \mathbb{N}$. 证明 $\cup A_n$ 是连通的.
\end{exercise}

\begin{proof}
    任取连续映射 \(  f: \bigcup A_{n}\to \left\{ 0,1 \right\}  \). 则 \(  f|A_{1}: A_1\to \left\{ 0,1 \right\}  \)也是连续的, 由于 \(  A_1  \)连通, \(  f|A_{1}  \)是常值的, 不妨设 \(  f\left( A_1 \right)= \left\{ 0 \right\}   \).  如果 \(  f\left( A_{n} \right)= \left\{ 0 \right\}   \), 则 \(  f\left( A_{n}\cap A_{n+ 1} \right)= \left\{ 0 \right\}   \), 由于 \(  A_{n+ 1}  \)是连通的,  \(  f|A_{n+ 1}  \)是常值映射, 从而 \(  f\left( A_{n+ 1} \right)= \left\{ 0 \right\}   \). 归纳可知 \[
    f\left( A_{n} \right)=\left\{ 0 \right\},\forall n \in \mathbb{N}  
    \]因此 \[
    f\left( \bigcup A_{n} \right)= \left\{ 0 \right\} 
    \] \(  f  \)是 \(  \bigcup A_{n}  \)上的常值映射. 这表明 \(  \bigcup A_{n}  \)是连通的.            
\end{proof}
\begin{exercise}{}{}
设 $A$ 是 $X$ 的连通子集. 那么 $A^\circ$ 和 $\partial A$ 是否一定是连通的? 反过来, 若 $A^\circ$ 和 $\partial A$ 都是连通的, $A$ 是否一定连通?
\end{exercise}
\begin{proof}
    \begin{enumerate}
        \item 若 \(  A  \)是 \(  X  \)的连通子集,  \(  A^{\circ}  \)和 \(   \partial A  \)均未必连通. 对于 \(  A^{\circ}  \), 考虑两个闭圆盘 \(  D^{2}  \)以边界\(   \partial D^{2}  \)上一点 \(  x_0  \)为基点的   一点并(Wedge Sum) \[
    A\coloneqq     D^{2}\vee_{x_0} D^{2}
        \] 由于 \(  D^{2}  \)是道路连通的, 其上每一点都有连接至 \(  x_0  \)的道路, 因此 \(  D^{2}\vee _{x_0}D^{2}  \)是道路连通的.       它的内部为 \[
        \left( D^{2}\vee_{x_0} D^{2} \right)^{\circ}=  \left( D^{2} \right)^{\circ} \sqcup \left( D^{2} \right)^{\circ}  
        \]是不连通的. 对于 \(   \partial A \) , 考虑一个圆环 \[
    A\coloneqq     \left\{ \left( x,y \right)\in \mathbb{R} ^{2}: x^{2}+ y^{2}\le 4  \right\}\cap \left\{ \left( x,y \right)\in \mathbb{R} ^{2}: x^{2}+ y^{2}\ge 1  \right\}
        \]是连通的, 但是它的边界为 \[
        \left( 2S^{1} \right)\cup \left( S^{1} \right)  
        \]其中 \(  \left( 2S^{1} \right)\cap \left( S^{1} \right)= \varnothing    \),  因此 \(   \partial A  \)不连通.
        
        \item 考虑 \(  A= \mathbb{Q}   \), 则 \(  \mathbb{Q} ^{\circ}= \varnothing  \)是连通的,  \(   \partial A= \mathbb{R}   \)也是连通的. 但是 \(  A  \)不连通.    
    \end{enumerate}
    
\end{proof}
\begin{exercise}{}{}
设 $n > 1, X$ 是 $\mathbb{R}^n$ 的可数子集. 证明 $\mathbb{R}^n - X$ 是路径连通的.
\end{exercise}
\begin{note}
    想法是, 如果我们不能直接提供一个好的路径, 我们就提供大量的``可能好''的候选路径, 并说明由于我们给的候选资源足够多, 其中总有一个路径不是坏的.
\end{note}
\begin{proof}
 任取 \(  p,q\in \mathbb{R} \setminus X  \), \(  p\neq q  \). 定义 \[
 H_{p,q}\coloneqq \left\{ x\in \mathbb{R} ^{n}: \left( x-\frac{p+ q }{2 }  \right) \cdot \left( p-q \right)= 0  \right\}
 \]  是过 \(  p,q  \)中点, 且与 \(  p-q  \)正交的超平面. \(  H_{p,q}  \)   与 \(  \mathbb{R} ^{n-1}  \)同胚. 任取 \(  r \in H_{p,q}  \), 考虑路径 \[
  \gamma _{r}\left( t \right)\coloneqq \begin{cases} \left( 1-2t \right)p+ 2tr,& t \in \left[ 0,\frac{1 }{2 }  \right]\\ 
   \left( 1-\left( 2t-1 \right)  \right)r+ \left( 2t-1 \right)q,& t \in \left[ \frac{1 }{2 },1  \right]      \end{cases}  
 \]  则 \(   \gamma _{r}\left( t \right)   \)是 \(  \mathbb{R} ^{n}  \)上连接了 \(  p,q  \)的路径. 断言, 存在 \(  r \in H_{p,q}  \)    , 使得 \(   \gamma _{r}\left( t \right)   \)是 \(  \mathbb{R} ^{n}\setminus X  \)上的道路.  若不然, 注意到对于不同的 \(  r_1,r_2\in H_{p,q}  \), \(   \gamma _{r_1},  \gamma _{r_2}  \)的像不交于 \(  p,q  \)以外 . 并且对于任意的 \(  r  \), 存在 \(  x_{r} \in X  \)使得 \(  x_{r} \in \operatorname{Im} \gamma _{r}  \). 设 $n > 1, X$ 是 $\mathbb{R}^n$ 的可数子集. 由选择公理, 存在单射 \[
\begin{aligned}
f: H_{p,q}&\to X\\ 
 r&\mapsto  x_{r}
\end{aligned}
 \]     但是 \(  \left| H_{p,q} \right|= \left| \mathbb{R} ^{n} \right|    \),   \(  H_{p,q}  \)是不可数集   , 矛盾. 因此存在 \(  r \in H_{p,q}  \), 使得 \(   \gamma _{r}  \)是 \(  \mathbb{R} ^{n}\setminus X  \)连接了 \(  p,q  \)的路径.    
\end{proof}
\begin{exercise}{}{}
若 $X$ 只有有限多个连通分支, 证明每个连通分支都是开集.
\end{exercise}

\begin{proof}
    设\(  C_{1},\cdots ,C_{n}  \)是 \(  X  \)的全体连通分支.  由于任意连通子集的闭包也是连通的,对于任意的 \(  i=  1,\cdots,n   \)  .\(  \overline{C_{i}}\supseteq C_{i}  \)是连通的, 而 \(  C_{i}  \)是极大的连通子集, 因此 \(  \overline{C_{i}}= C_{i}  \), \(  C_{i}  \)是一个闭集.  又\(  C_1,\cdots ,C_{n}  \)两两无交, 我们有对于任意的 \(  i=  1,\cdots,n   \), \[
    C_{i}= X\setminus \left( \bigcup _{j\neq i}C_{j} \right) 
    \]是一个开集.  
\end{proof}
\begin{exercise}{}{}
设 $A, B$ 是 $X$ 的非空子集, 满足 $X = A^\circ \cup B^\circ$. 若 $X$ 是路径连通空间, 证明 $A$ 的任一路径连通分支与 $B$ 的交都是非空的.
\end{exercise}

\begin{proof}
  由于 \(  X  \)是道路连通的,  \(  X  \)也是连通的. 
  
  如果 \(  C  \)是 \(  A  \)的路径连通分支, 使得 \(  C\cap B= \varnothing  \), 则 \[
  C\subseteq X\setminus B\subseteq A^{\circ}
  \]
  若 \(  B^{\circ}= \varnothing  \), 则 \(  X= A^{\circ}  \), \(  A^{\circ}  \)是路径连通的, 从而 \(  C= A^{\circ}  \). 但是 \(  B  \)非空,  \(  B\cap C= B\cap X= B\neq \varnothing  \)矛盾.
  
  因此 \(  B^{\circ}\neq \varnothing  \), 取 \(  b \in B^{\circ}  \). 由于 \(  X  \)是道路连通的, 对于任意的 \(  x \in C  \), \(  x \in A^{\circ}  \), 存在 \(  x  \)到 \(  b  \)的路径 \(   \gamma   \).
  令 \[
  S\coloneqq \left\{ t\in \left[ 0,1 \right]:  \gamma \left( [0,t] \right)\subseteq A ^{\circ}  \right\}
  \]令 \[
  t_0= \sup S
  \]则 \(  t_0\in \left[ 0,1 \right]  \). 令 \(  p=  \gamma \left( t_0 \right)   \). 断言 \(  t_0<  1  \), 否则 \(   \gamma \left( [0,1) \right)\subseteq A^{\circ}   \), 注意到 \(   \gamma \left( [0,1) \right)\cup C   \)是道路连通的, 从而 \(   \gamma \left( [0,1) \right)\subseteq C   \)  , \(   \gamma \left( 1 \right)\in B^{\circ}   \) . 由于 \(  B^{\circ}  \)是开集, 存在 \(   \gamma \left( 1 \right)   \)在 \(  B^{\circ}  \)中的开邻域 \(  W  \), 使得 \(   \gamma ^{-1} \left( W \right)   \)是一个开集, 从而存在 \(  t_1< 1  \), 使得 \(  t_1 \in  \gamma ^{-1} \left( W \right)   \),  \(   \gamma \left( t_1 \right)\in W\subseteq B^{\circ}   \)    . 这表明 \(   \gamma \left( t_1 \right)\in C\cap B   \), 与 \(  C\cap B= \varnothing  \)矛盾.     
  \begin{enumerate}
    \item 若 \(  p \in A^{\circ}  \), 由于 \(  A^{\circ}  \)是开集,存在 \(  p  \)在 \(  A^{\circ}  \)中的开邻域 \(  U  \), 由于 \(   \gamma   \)连续, \(   \gamma ^{-1} \left( U \right)   \)是一个开集. 又 \(  t_0< 1  \), 存在开区间 \(  \left( c,d \right)   \)使得 \( t_0\in \left( c,d\right)\subseteq  \gamma ^{-1} \left( U \right)   \). 则          \(  \frac{t_0+ d }{ 2}> t_0   \), \(  \frac{t_0+ d }{2 }\in S   \), 与 \(  t_0= \sup S  \)矛盾.
    \item 因此只能有 \(  p\in B^{\circ}  \), 但根据 \(  t_0  \)的定义,   \(   \gamma \left( [0,t_0) \right)\subseteq A^{\circ}   \),  \(   \gamma \left( t_0 \right)\in B^{\circ}   \).而注意到 \[
     \gamma \left( [0,t_0) \right)\cup C 
    \]是道路连通的, 因此 \(   \gamma \left( [0,t_0) \right)\subseteq C   \). 由于 \(  B^{\circ}  \)是开集, 存在 \(  p  \)在 \(  B^{\circ}  \)中的开邻域 \(  V  \), 使得 \(   \gamma ^{-1} \left( V \right)   \)是包含了 \(  t_0  \)的一个开集. 进而存在开区间 \(  \left( e,f \right)   \), 使得 \(  t_0\in \left( e,f \right)\subseteq  \gamma ^{-1} \left( V \right)    \), 从而 \(   \gamma \left( \left( e,f \right)  \right)\subseteq B^{\circ}   \)         .此时 \(  \gamma  \left( \frac{e+ t_0 }{2 }   \right) \in  \gamma \left( [0,t_0) \right)\cap B^{\circ}\subseteq C\cap B   \) .与 \(  C\cap B= \varnothing  \)矛盾. 
  \end{enumerate}
综上所述,  不存在 \(  A  \)的连通分支 \(  C  \), 使得 \(  C\cap B= \varnothing  \).   
\end{proof}
\begin{exercise}{}{}
设 $X$ 是局部路径连通空间. 证明 $X$ 的任一连通开集都是路径连通的.
\end{exercise}
\begin{proof}
    任取连通开集 \(  U  \). 取定 \(  x_0 \in U  \), 定义 \[
    W\coloneqq \left\{ x\in U: \text{存在}U\text{上连接}x_0\text{和}x\text{的道路} \right\}
    \]  若 \(  x \in W  \), 由于 \(  X  \)是局部道路连通的, 存在 \(  x  \)在 \(  U  \)中的道路连通的开邻域 \(  V  \). \(  V  \)上任一点存在连接到 \(  x  \)的道路, 又 \(  x \in W  \), 因此 \(  V  \)上任一点存在连接到 \(  x_0  \)的道路     . 这表明 \(  V  \)是 \(  x  \)在 \(  U \)上的开邻域, 使得 \(  x \in V\subseteq U  \) .   因此 \(  W  \)是\(  U  \)上的 开集.      
    
    另一方面, 任取 \(  y \in U\setminus W  \), 存在 \(  y  \)在 \(  U  \)中的道路连通的开邻域 \(  V  \). 任取 \(  p\in V  \), 存在\(  U  \)上连接了 \(  p  \)和 \(  y  \)的道路. 如果 \(   p \in W  \), 则    存在 \(  U  \)上连接了 \(  p  \)和 \(  x_0  \)的道路,           进而存在 \(  U  \)上连接 \(  y  \)和\(  x_0  \)的道路, 矛盾. 因此 \(  p\not \in W  \), \(  p \in U\setminus W  \). 这表明 \(  V  \)是 \(  x \)在 \(  U \)上的一个开邻域, 使得 \(  x \in V\subseteq U\setminus W  \) . 因此 \(  U\setminus W  \)是一个开集. 现在我们有 \[
    U= W\cup \left( U\setminus W \right) 
    \]是两个开集的并, 由于 \(  U  \)是连通的, 并且 \(  x_0 \in W  \), \(  W  \)非空, 只能有 \(  W= U  \). 因此 \(  U  \)是道路连通的.                        
\end{proof}
\begin{exercise}{}{}
 设 $X$ 是 $\mathbb{R}^2$ 中在线段 $[0,1] \times 0$ 上的有理点, $T$ 是从点 $p=(0,1)$ 到 $X$ 中的点的线段之并集.
\begin{enumerate}
    \item 证明 $T$ 是路径连通的, 但却只在点 $p$ 处是局部连通的.
    \item 构造 $\mathbb{R}^2$ 的一个子集, 它是路径连通的, 但在每一点处都不是局部连通的.
\end{enumerate}
\end{exercise}
\begin{proof}
    \begin{enumerate}
        \item 任取 \(  T  \)上两点 \(  q_1,q_2  \), 则 \(  q_1,q_2  \)落在与 \(  p  \)相连的某个线段上. 从而存在 \(  q_1  \)到 \(  p  \)的路径 \(   \gamma _1   \)和 \(  p  \)到 \(  q_2  \)的路径 \(   \gamma _2   \). 因此 \[
         \gamma \coloneqq \begin{cases}  \gamma _1 \left( 2t \right),& t\in \left[ 0,\frac{1 }{2 }  \right]\\ 
           \gamma _2 \left( 2t-1 \right),& t\in \left[ \frac{1 }{2 },1  \right]     \end{cases} 
        \]就是 \(  T  \)上连接了 \(  q_1  \)和 \(  q_2  \)的路径. 因此 \(  T  \)是路径连通的.
        
        任取以 \(  p  \)为中心的开球 \(  B  \),  \(  B  \)上任一点都落在与 \(  p  \)相连的某个线段上, 由于以上关于 \(  T  \)的论证是基于\(  p  \)点的,   与之类似的过程可以给出 \(  B  \cap T\)是路径连通的, 进而是连通的.
        
        但是对于任意 \(  T  \)上异于 \(  p  \)的一点 \(  q  \), 令 \(  d = d\left( p,q \right)   > 0\), 考虑开球 \(  B_{d}\left( q \right)   \). 考虑连续映射 \(   \pi:T\setminus \left\{ p \right\}\to \mathbb{Q} \cap \left[ 0,1 \right]   \), \(   \pi\left( r \right)   \)为 \(  r  \)所在线段对应的有理数.   任取 \(  B_{d}\left( q \right)   \)的非空开子集 \(  U  \),    任取 \(  q_0\in U\cap T  \). 不难发现\(  U  \)与无穷多个线段相交. 设 \(  C_{q_0}  \)是 \(  q_0  \)所在的\(  U  \)的 连通分支. 由于 \(  \mathbb{Q}   \)是完全不连通的,  \(   \pi\left( C_{q_0} \right)   \)只能是单点集.     但是上面的事实表明 \(   \pi\left( U\cap T \right)   \) 是无限集, 这表明 \(  C_{q_0}\neq U \cap T \), \(  U\cap T  \)不是连通的. 因此 任意 异于 \(  p  \)的点 \(  q  \)处都不是局部连通的.
        
        \item 令 \(  X= \left( \left[ 0,\frac{1 }{2 }  \right]\cap \mathbb{Q}   \right)\times \left\{ 0 \right\}   \), \( Y= \left( \left[ \frac{1 }{2 } ,1 \right]\cap \mathbb{Q} \right)  \times \left\{ 1 \right\}    \). 令 \(  p_1= \left( \frac{1 }{2 },1  \right)   \), \(  p_2= \left( \frac{1 }{2 },0  \right)   \).   令 \(  T_1  \)是点 \(  p_1  \)到 \(  X  \)中的线段之并, \(  T_2  \)是点 \(  p_2  \)到 \(  Y  \)中的线段之并. 则 \[
        T_1\cap T_2= \left\{ \frac{1 }{2 }  \right\}\times \left[ 0,1 \right] 
        \]       根据1. 的讨论, 容易得到 \(  T_1  \)是道路连通的, 且在  \(  T_1\setminus T_2  \)之外任一点处都不是局部连通的.  类似地 \(  T_2  \)是道路连通的, 且在 \(  T_1\setminus T_2  \)外任一点处都不是局部连通的. 又 \(  T_1\cap T_2  \)是道路连通的, 因此 \(  T_1\cup T_2  \)是道路连通的, 要说明 \(  T_1\cup T_2  \)在每一点处都不是局部连通的, 只需要说明 \(  T_1\cup T_2  \)在 \(  T_1\cap T_2  \)上任一点处都不是局部连通的. 具体地, 任取 \(  q \in T_1\cap T_2  \),令 \[
        d=  \max \left\{ d\left( q,p_1 \right),d\left( q,p_2 \right)   \right\}
        \]则 \( \left\{ p_1,p_2 \right\}\not \subseteq  B_{d}\left( q \right) \) .  不妨设 \(  p_1\not \in B_{d}\left( q \right)   \)  , 任取 \(  B_{d}\left( q \right)   \)的非空开子集 \(  U  \),            任取 \(  q_0\in U\cap T_1  \), \(  U  \)与 \(  T_1  \)中的无穷多个线段相交. 类似地, 考虑映射 \(   \pi:T_1\cup T_2\setminus \left\{ p_1,p_2 \right\}\to \mathbb{Q} \cap \left[ 0,1 \right]   \) ,  \(   \pi\left( r \right)   \)为 \(  r  \)所在线段对应的有理数.   设 \(  C_{q_0}  \)是 \(  q_0  \)所在的 \(  U  \)的连通分支, 由于 \(  \mathbb{Q}   \)完全不连通, \(   \pi\left( C_{q_0} \right)   \)是单点集, 但是 \(   \pi\left( U\cap T_1 \right)   \)是无限集, 这表明 \(  C_{q_0}\neq U  \cap T_1\), \(  U \cap T_1 \)不是连通的. 这表明任取 \(  q \in T_1\cap T_2  \)处都不是局部连通的.            
    \end{enumerate}
    
\end{proof}
\begin{exercise}{}{}
    设 $X \neq \emptyset$ 是拓扑空间. 我们称 $X$ 是\dfntxt{不可约的}, 是指对 $X$ 中任何闭集 $P, Q$ 使得 $X = P \cup Q$, 则必有 $X = P$ 或 $X = Q$.
    
 
\begin{enumerate}
    \item 证明若 $X$ 是不可约的, 则任一非空开集 $U \subset X$ 也是不可约的, 并且在 $X$ 中稠密.
    \item 设 $X$ 为拓扑空间, $A \subset X$ 是不可约子空间. 证明闭包 $\overline{A} \subset X$ 也是不可约的.
    \item 设 $X$ 是不可约的拓扑空间. 证明 $X$ 是连通的.
    \item 一个空间能同时是不可约的而且 Hausdorff 的吗? 能的话都找出来.
    \item 设 $X, Y$ 为拓扑空间, $A \subset X$ 是不可约子空间. 证明对于任意连续映射 $f: X \to Y, f(A) \subset Y$ 也是不可约的.
\end{enumerate}
\end{exercise}
\begin{proof}
    \begin{enumerate}
        \item 任取 \(  U  \)中的闭集 \(  P\cap U  \), \(  Q\cap U  \), 其中 \(  P,Q  \)是 \(  X  \)中的闭集, 如果 \[
        U= \left( P\cap U \right)\cup \left( Q\cap U \right)  
        \]则 \[
        X= \left( P\cup Q \right)\cup \left( X\setminus U   \right) 
        \]     是两个闭集的并, 一定有 \(  X= P\cup Q  \)或 \(  X= X\setminus U   \)  , 对于前者, 必有 \(  X= P  \)或 \(  X= Q  \), 此时 \(  U= P\cap U  \) 或\(  U= Q\cap U  \). 对于后者, 只能有 \(  U= \varnothing \), 但是 \(  U  \)非空, 矛盾. 这表明    \(  U  \)是不可约的. 此外, 注意到 \[
        X= \left( X\setminus U \right)\cup \bar{U} 
        \]   是两个闭集的并, 由于\(  X  \)不可约, 且  \(  U  \)非空, 只能有 \(  X= \bar{U}  \) , 故 \(  U  \)在 \(  X  \)中稠密.   
        \item 任取 \(  \bar{A}  \)中的闭集 \(  P\cap \bar{A}  \)  , \(  Q\cap \bar{A}  \), 其中 \(  P,Q  \)是 \(  X  \)中的闭集. 使得 \[
        \bar{A}= \left( P\cap \bar{A} \right)\cup \left( Q\cap \bar{A} \right)  
        \]我们有 \[
        A=  \left( P\cup Q \right)\cap A= \left( P\cap A \right)\cup \left( Q\cap A \right)  
        \]只能有 \(  A= Q\cap A  \)或 \(  A= P\cap A  \)  , 若 \(    A= Q\cap A\) 则 \[
        \bar{A}= \overline{Q\cap A}\subseteq  \overline{Q}= Q\implies  \bar{A}=  Q\cap \bar{A}
        \]   类似地, 若 \(  A= P\cap A  \), 则 \(  \bar{A}=  P\cap \bar{A}  \). 因此 \(  \bar{A}  \)也是不可约的.
        \item 如果 \(  X  \)是不连通的, 则存在开集 \(  U,V  \), 使得 \[
        X= U\cup V,\quad U\cap V= \varnothing,\quad U,V\neq X
        \]  此时注意到 \(  U= X\setminus V  \)和 \(  V= X\setminus U  \)都是闭集, 只能有 \(  X= U  \)或\(  X= V  \), 矛盾.  因此 \(  X  \)是连通的.     
        \item 如果 \(  X  \)是Hausdorff的, 且不是单点的空间. 任取 \(  x,y \in X  \), \(  x \neq y  \). 存在 \(  x  \)的开邻域 \(  U  \)和 \(  y  \)的开邻域 \(  V  \)使得 \(  U\cap V= \varnothing  \).则 \[
        X=X\setminus \left( U\cap V \right)=   \left( X\setminus U \right)\cup \left( X\setminus V \right)  
        \]  由于 \(  X  \)是不可约的, 要么 \(  X= X\setminus U  \), 要么 \(  X= X\setminus V  \). 则 \(  U= \varnothing  \)或 \(  V= \varnothing  \), 与 \(  x \in U  \)或 \(  y \in V  \)矛盾. 因此不存在具有两个点以上的不可约的Hausdorff空间.   这样  空间在同胚的意义下有且仅有 一种\[
     \left( X,\tau  \right)= \left( \left\{ x \right\}, \left\{ \varnothing, \left\{ x \right\} \right\} \right)  
        \]
        \item 任取 \(  f\left( A \right)   \)上的闭集 \(  P\cap f\left( A \right), Q\cap f\left( A \right)    \), 其中 \(  P,Q  \)是 \(  Y  \)上的闭集, 使得 \[
        f\left( A \right)= \left( P\cap f\left( A \right)  \right)\cup \left( Q\cap f\left( A \right)  \right)   
        \]   则 \[
        A\subseteq f^{-1} \left( P \right)\cup f^{-1} \left( Q \right)  
        \]         \[
        A= \left( f^{-1} \left( P \right)\cap A  \right)\cup \left( f^{-1} \left( Q \right)\cap A  \right)  
        \]    其中 \(  f^{-1} \left( P \right), f^{-1} \left( Q \right)    \)是 \(  X  \)上的闭集. 由于 \(  A  \)是不可约的, 要么 \(  A\subseteq f^{-1} \left( P \right)   \), 要么 \(  A\subseteq f^{-1} \left( Q \right)   \). 从而 \(  f\left( A \right)\subseteq P   \)或 \(  f\left( A \right)\subseteq Q   \). 这表明 \(  f\left( A \right)   \)是不可约的.        
    \end{enumerate}
    
\end{proof}
\end{document}